%==============================================================================
% Sjabloon poster bachproef
%==============================================================================
% Gebaseerd op document class `a0poster' door Gerlinde Kettl en Matthias Weiser
% Aangepast voor gebruik aan HOGENT door Jens Buysse en Bert Van Vreckem

\documentclass[a0,portrait]{hogent-poster}

% Info over de opleiding
\course{Bachelorproef}
\studyprogramme{toegepaste informatica}
\academicyear{2025-2026}
\institution{Hogeschool Gent, Valentin Vaerwyckweg 1, 9000 Gent}

% Info over de bachelorproef
\title{Een machine learning benadering voor het identificeren van relevante fysieke spelersattributen in voetbal scouting}
\author{Vanacker Matthijs}
\email{matthijs.vanacker@student.hogent.be}
\supervisor{Dhr. T. Antjon}
\cosupervisor{Dhr. A. Martiny (KAA Gent)}

% Indien ingevuld, wordt deze informatie toegevoegd aan het einde van de
% abstract. Zet in commentaar als je dit niet wilt.
\specialisation{AI \& Data Engineering}
\keywords{Sport, Scouting, Machine Learning, Voetbal}
\projectrepo{https://github.com/MatthijsVanacker/latex-hogent-bachproef-25-26}

\begin{document}

\maketitle

\begin{abstract}
De hoeveelheid data die beschikbaar is in de moderne voetbalwereld blijft groeien. Voetbalclubs zoeken voortdurend naar manieren om deze data toe te passen en zo vooruitgang te boeken in alle delen van de club, van de eerste ploeg tot de jeugdacademie. Een sector waar de groeiende hoeveelheid data een grote impact heeft, is het scoutingsdepartement. 
In dit onderzoek wordt gebruikgemaakt van machine learning-technieken om belangrijke fysieke spelersattributen te identificeren en deze inzichten te gebruiken om spelers een fysieke score toe te kennen. Deze aanpak biedt scouts en data-analisten een alternatief voor de huidige berekening van de fysieke score van spelers, gebruikmakend van een experteninschatting van het belang van fysieke attributen. 
De centrale onderzoeksvraag van dit onderzoek luidt: \emph{Hoe kan een objectieve en reproduceerbare fysieke spelersscore ontwikkeld worden op basis van fysieke wedstrijddata met behulp van machine learning-technieken?} 
Om dit te onderzoeken worden fysieke wedstrijddata gekoppeld aan wedstrijdresultaten over 25 competities, verkregen via de SkillCorner API. Na de data op te schonen, feature engineering toe te passen en te normaliseren per positiegroep en competitie worden regressie- en classificatiemodellen getest op drie datasets: een \emph{individuele}, een \emph{geaggregeerde} en een \emph{verschil}-dataset. Het best presterende model is een XGBoost-classificatiemodel getraind op de geaggregeerde dataset met een nauwkeurigheid van 44\%. Uit dit model worden de belangrijkste fysieke attributen geïdentificeerd via SHAP-waarden, deze worden vervolgens gebruikt om voor elke speler, per competitie en per seizoen, een fysieke score te berekenen tussen 0 en 100. Om een overzicht te krijgen van de spelers en hun score wordt deze score gevisualiseerd via een extra pagina binnen het bestaande fysieke scoutingsdashboard van de club.
De resultaten van dit onderzoek tonen aan dat fysieke wedstrijddata alleen slechts een beperkte voorspellende waarde heeft voor wedstrijdresultaten. Desondanks is er aangetoond dat het mogelijk is een objectieve en reproduceerbare fysieke spelersscore te creëren op basis van machine learning-technieken. 

\end{abstract}

\begin{multicols}{2} % This is how many columns your poster will be broken into, a portrait poster is generally split into 2 columns

\section{Introductie}

De club krijgt zeer veel fysieke wedstrijddata binnen via SkillCorner, het exacte belang van iedere fysieke spelersstatistiek om een positief resultaat te behalen, is echter moeilijk te beoordelen. De huidige aanpak is gebaseerd op de inschatting van experten op de invloed van enkele cruciale fysieke attributen. Om optimaal gebruik te maken van al deze gegevens kan machine learning ingezet worden, zo'n algoritme kan al deze fysieke data analyseren en de belangrijkste attributen identificeren die bijdragen aan het behalen van een overwinning. 
\newline \newline
Hieruit volgt de \textbf{centrale onderzoeksvraag}:
\newline \newline
\emph{Hoe kan een objectieve en reproduceerbare fysieke spelersscore ontwikkeld worden op basis van fysieke wedstrijddata met behulp van machine learning-technieken?} 

\section{Data en modeltraining}
\subsection{Verschillende datasets}
Alle kwantitatieve fysieke wedstrijddata wordt omgezet naar relatieve waarden per 90 minuten. De focus ligt op spelers die minstens 75 minuten gespeeld hebben in een wedstrijd, waarna de data wordt genormaliseerd per positiegroep en competitie. Er worden drie verschillende datasets gecreëerd:
\begin{itemize}
  \item \textbf{Individuele dataset}: Deze dataset bevat individuele fysieke wedstrijddata van spelers.
  \item \textbf{Geaggregeerde dataset}: Deze dataset bevat geaggregeerde fysieke wedstrijddata van spelers per club en wedstrijd, gegroepeerd per positiegroep.
  \item \textbf{Verschil-dataset}: Deze dataset bevat het verschil van geaggregeerde data van twee teams in een wedstrijd.
\end{itemize}
\subsection{Soorten modellen}
\begin{itemize}
  \item \textbf{Classificatiemodel}: Dit model voorspelt het resultaat van een wedstrijd gebaseerd op de prestaties van een team of speler.
  \item \textbf{Regressiemodel}: Dit model voorspelt een getal, zoals het aantal gescoorde doelpunten, tegendoelpunten of het doelpuntenverschil van een wedstrijd.
\begin{itemize}
  \item \textbf{Hybride aanpak}: Binnen de regressiemodellen wordt ook een hybride aanpak gehanteerd, waarbij het voorspelde doelpuntenverschil vertaald wordt naar een voorspelling van het wedstrijdresultaat.
\end{itemize}
\end{itemize}

\subsection{Gekozen combinatie}
Na het evalueren van de verschillende modellen in combinatie met de verschillende datasets verkreeg een XGBoost-classificatiemodel getraind op de geaggregeerde dataset, met een nauwkeurigheid van 44\% de beste score.
\section{Berekening van de fysieke score}
\begin{enumerate}
  \item Na het trainen van de modellen worden SHAP-waarden (SHapley Additive exPlanations) gebruikt om de bijdrage van elke \emph{feature} aan de voorspelling van het model te bepalen.
  \item Deze bijdragen worden gegroepeerd per positie en omgezet naar procentuele waarden.
  \item De berekende waarden worden vervolgens gebruikt om een fysieke score te berekenen voor elke speler per wedstrijd, afhankelijk van zijn positie.
  \item Hierna wordt een gewogen gemiddelde van deze scores genomen over alle wedstrijden van een speler in een competitie en seizoen.
  \item Tot slot worden deze scores omgezet naar een schaal van 0 tot 100.
\end{enumerate}
\section{Dashboard implementatie}
De berekende fysieke scores worden gevisualiseerd in een toegevoegde pagina binnen het bestaande fysieke scoutingsdashboard van de club.
\begin{center}
  \captionsetup{type=figure}
  \includegraphics[width=1.0\linewidth]{dashboard_bovenkant}
  \captionof{figure}{De machine learning-pagina van het scoutingdashboard}
\end{center}

\section{Conclusies}

Dit onderzoek toont aan dat het mogelijk is om een objectieve, reproduceerbare fysieke spelersscore te creëren met machine learning-technieken. Desondanks toont de lage nauwkeurigheid van het model aan dat een wedstrijdresultaat voorspellen alleen op basis van fysieke spelersdata een grote uitdaging blijft.

\section{Toekomstig onderzoek}

Tijdens dit onderzoek zijn er mogelijkheden voor alternatief onderzoke naar voren gekomen, de voornaamste hiervan zijn:
\begin{itemize}
  \item Kan een clusteringmodel ingezet worden om fysieke \emph{spelerstypes} te identificeren en kunnen \emph{elite} spelers binnen deze types worden gevonden?
  \item Voortbouwend op het voorgaande, wat is het effect van meerdere van deze \emph{elite} spelers in een wedstrijd op het resultaat?
\end{itemize}

\end{multicols}
\end{document}